% Zone „Das kennst du schon“ – Kreis, Kl. 8
\setcounter{aufgabe}{0}
\einheitenkopf*{\hypertarget{zone}{}Das kennst du schon}

\begin{aufgabe}{Ich kann mit Kommazahlen multiplizieren und dividieren und das Ergebnis runden. Berechne.}
\begin{teilezwei}
\tz $3 \cdot 1{,}2 =$ \leerfeld & \tz $0{,}5 \cdot 8 =$ \leerfeld \\
\tz $2{,}5 \cdot 1{,}4 =$ \leerfeld & \tz $7{,}5 : 2{,}5 =$ \leerfeld \\
\end{teilezwei}
Runde auf die angegebene Stelle.
\begin{teilezwei}
\tz $12{,}46$ auf eine Stelle: \leerfeld & \tz $8{,}0472$ auf zwei Stellen: \leerfeld \\
\end{teilezwei}
\begin{teile}
\teil Rechne $1{,}26 \cdot 5$ und runde erst das Ergebnis auf eine Stelle. Vergleiche mit der Rechnung, bei der du zuerst $1{,}26$ auf $1{,}3$ rundest. Welches Ergebnis stimmt?
\end{teile}
\end{aufgabe}

\begin{aufgabe}{Ich kann Umfang und Flächeninhalt unterscheiden und die richtige Einheit angeben. Ein Rechteck ist 3\,cm lang und 2\,cm breit, ein Quadrat hat die Seitenlänge 1,5\,m. Berechne mit Einheit.}
\begin{teilezwei}
\tz Rechteck, Umfang: \leerfeld & \tz Rechteck, Flächeninhalt: \leerfeld \\
\tz Quadrat, Umfang: \leerfeld & \tz Quadrat, Flächeninhalt: \leerfeld \\
\end{teilezwei}
Kreuze die passende Einheit an.
\begin{teile}
\teil Länge des Zauns um ein Gemüsebeet: 26 \quad \kreuz{m}\quad\kreuz{m²}
\teil Größe der Rasenfläche in einem Garten: 40 \quad \kreuz{m}\quad\kreuz{m²}
\end{teile}
\end{aufgabe}

\begin{aufgabe}{Ich kann eine Formel nach einer anderen Größe umstellen. Stelle nach der Größe in Klammern um.}
\begin{gleichungsraster}[1]
\gl{A = a \cdot h \quad (h)} & \gl{s = v \cdot t \quad (t)} \\
\end{gleichungsraster}
\medskip
Löse nach der Variablen auf.
\begin{gleichungsraster}[2]
\gl{7{,}5 = 2{,}5 \cdot x} & \gl{60 = 2 \cdot 3 \cdot z} \\
\gl{49 = a^2} & \gl[3]{50 = 2 \cdot a^2} \\
\end{gleichungsraster}
\end{aufgabe}

\begin{aufgabe}{Ich kann Zahlen quadrieren und die Quadratwurzel ziehen. Berechne.}
\begin{teilezwei}
\tz $6^2 =$ \leerfeld & \tz $\sqrt{81} =$ \leerfeld \\
\tz $1{,}2^2 =$ \leerfeld & \tz $0{,}3^2 =$ \leerfeld \\
\tz $\sqrt{2{,}25} =$ \leerfeld & \tz $\sqrt{40} \approx$ \leerfeld\ (eine Stelle) \\
\end{teilezwei}
\end{aufgabe}

\begin{aufgabe}{Ich finde den Fehler beim Quadrieren. Paul berechnet den Flächeninhalt eines Quadrats mit der Seitenlänge $a = 8$\,cm so:}
\rechnung{A &= a^2 \\ A &= 2 \cdot 8\,\text{cm} \\ A &= 16\,\text{cm}^2}
\begin{teile}
\teil Finde den Fehler, benenne ihn und korrigiere die Rechnung.
\teil Berechne den Flächeninhalt eines Quadrats mit der Seitenlänge $a = 1{,}1$\,m.
\end{teile}
\end{aufgabe}

\begin{aufgabe}{Ich kann mit rechtem, gestrecktem und vollem Winkel rechnen. Gib die Größe in Grad an.}
\begin{teilezwei}
\tz Vollwinkel: \leerfeld[°] & \tz rechter Winkel: \leerfeld[°] \\
\tz gestreckter Winkel: \leerfeld[°] & \tz $360° : 8 =$ \leerfeld[°] \\
\end{teilezwei}
Ergänze.
\begin{teilezwei}
\tz $110° +$ \leerfeld[°] $= 360°$ & \tz $65° +$ \leerfeld[°] $= 180°$ \\
\end{teilezwei}
\end{aufgabe}

\begin{aufgabe}{Ich kann einen Anteil als Bruch und in Prozent angeben. Gib den Anteil als gekürzten Bruch und in Prozent an.}
\sachtabelle{lcc}{Anteil & gekürzter Bruch & in Prozent}{a) 3 von 12 & \leerzelle & \leerzelle \\ b) 7 von 14 & \leerzelle & \leerzelle \\ c) 18 von 24 & \leerzelle & \leerzelle \\ d) 45 von 360 & \leerzelle & \leerzelle \\ e) 7 von 30 (Prozent auf eine Stelle) & \leerzelle & \leerzelle}
\setcounter{teil}{5}
\begin{teile}
\teil Von 40 Äpfeln in einer Kiste sind 10 faul. Welcher Anteil der Äpfel ist nicht faul? \\ Bruch: \leerfeld \quad Prozent: \leerfeld[\%]
\end{teile}
\end{aufgabe}
